\input{../universal-pre}

\title{\tbf{AB1995}}
\author{\tbf{Notes} by Cong Wen, 03/2022}
\date{}

\begin{document}
\maketitle
\thispagestyle{empty}
\begin{abstract}
  This is my self-use notes\footnote{Last updated on March 08, 2022} on \tbf{Hilbert-Samuel theorem}\cite{AB1995}. Please reach out if you find anything incorrect.
\end{abstract}
\tableofcontents


% ----------------------------------------------------------------------------------------------------


\section{introduction}
\begin{stp}
  $K$ number field, $\cO_{K}$ its ring of integers. $\Phi$ embeddings of $K$ to $\bC$ modulo complex conjugation. $X\ar\Spec\cO_{K}$ flat projective scheme, with smooth fiber of dimension $r\geq 1$. $\cL$ ample line bundle over $X$, for each place $\sigma\in\Phi$, endow $\cL_{\sigma}, \sigma\in\Phi$ with Hermitian metric with positive curvature.


  There is a canonical embedding, making the free $\cO_{K}$-module $\HH{}{0}(X,n\cL)$ into a lattice in $\bR$-vector space $W_{n}=\HH{}{0}(X,n\cL)\otimes_{\bZ}\bR$
  \[
    0\ar\HH{}{0}(X,n\cL)\ar\bigoplus_{\sigma\in\Phi}\HH{}{0}(X,n\cL)\otimes K_{\sigma}=W_{n}.
  \]
  There are two possible norms on $W_{n}$,
  \begin{enumerate}
    \item (Sup norm), for any $x=(x_{\sigma})_{\sigma\in\Phi}\in W_{n}$, $B_{n}\subset W_{n}$ the unit ball.
          \[
            \dP{x}_{\mathrm{sup}}=\sup_{\sigma\in\Phi}\sR{\sup_{p\in X_{\sigma}}\sP{s(p)}}.
          \]
    \item ($L^{2}$ norm), for any $x=(x_{\sigma})_{\sigma\in\Phi}\in W_{n}$, $C_{n}\subset W_{n}$ the unit ball.
          \[
            \dP{x}_{\mathrm{sup}}=\sup_{\sigma\in\Phi}\sR{\int_{X_{\sigma}}\sP{s(p)}^{2}\dif p}.
          \]
  \end{enumerate}

  The \tbf{Arakelov sections} is defined
  \[
    \HH{\mathrm{Ar}}{0}(X,n\cL)=\HH{}{0}(X,n\cL)\cap B_{n}.
  \]
\end{stp}
\begin{thm}[Sup unit ball]\label{thm:supub}
  The following limit exists, and is denoted by $(\cL)^{r+1}$,
  \[
    (\cL)^{r+1}=\lim_{n\ar\infty}\frac{(r+1)!}{n^{r+1}}\log\frac{\Vol(B_{n})}{\Vol(\HH{}{0}(X,n\cL))}.
  \]
\end{thm}

\begin{thm}[$L^{2}$ unit ball]\label{thm:L2ub}
  The following limits exists, and is equal to $(\cL)^{r+1}$ above,
 \[
   (\cL)^{r+1}=\lim_{n\ar\infty}\frac{(r+1)!}{n^{r+1}}\log\frac{\Vol(C_{n})}{\Vol(\HH{}{0}(X,n\cL))}.
  \]
\end{thm}

\begin{cor}
  If $(\cL)^{r+1}>0$, then $\HH{\mathrm{Ar}}{0}(X,n\cL)\neq 0, n\gg 0$.
\end{cor}


\section{additive volumes}
\begin{stp}
  Let $M$ be a finite type $\cO_{K}$-module, choose a volume form $\eta_{\sigma}$ on $M\otimes_{\sigma}\bC$ for each $\sigma\in\Phi$. This gives a volume form $\eta=\bigotimes_{\sigma\in\Phi}\eta_{\sigma}$ on $M_{\bR}=M\otimes_{\bZ}\bR=\bigoplus_{\sigma\in\Phi}M\otimes_{\sigma}\bC$. Note $\oL{M}=M/M_{\mathrm{tor}}\ahr M_{\bR}$ and that $\cO_{K}$ has a canonical volume form.

\end{stp}

\begin{dfn}
  Define $\chi(M,\eta)$ by
  \[
    \chi(M,\eta)=\log\frac{\sP{M_{\mathrm{tor}}}}{\Vol(M_{\bR}/\oL{M})}.
  \]
\end{dfn}
\begin{prp}
  Take $(M^{i},\eta^{i}), \eta^{i}\in\Det(M^{i}_{\bR}), i=1,2,3$ such that
  \[
    0\ar M^{1}\ar M^{2}\ar M^{3}\ar 0, \eta^{2}=\eta^{1}\otimes\eta^{3},
  \]
  then $\chi(M^{2},\eta^{2})=\chi(M^{1},\eta^{1})+\chi(M^{3},\eta^{3})$.
\end{prp}

\begin{dfn}\label{dfn:volk}
  Let $A=\oplus_{n\geq 0}A_{n}$ be a graded $\cO_{K}$-algebra of finite type, $M=\oplus_{n\geq 0}M_{n}$ be a graded $A$-module of finite type. Denote by $P$ the Hilbert polynomial of $M\otimes_{\cO_{K}}K$.

  For any real number $k$, define volume form $\alpha_{n,k}$ such that
  \[
    \chi(M_{n},\alpha_{n,k})=k\sum_{i=0}^{n-1}P(i)+\chi(\cO_{K})P(n).
  \]
\end{dfn}
\begin{prp}
  Let $k$ be any real number, $M^{i}, i=1,2,3$ are graded $A$-modules of finite type with
  \begin{align*}
    0\ar M^{1}\ar M^{2}\ar M^{3}\ar 0,\\
    \phi_{n}:\Det(M^{2}_{n}\otimes_{\bZ}\bR)\cong\Det(M^{1}_{n}\otimes_{\bZ}\bR)\otimes\Det(M^{3}_{n}\otimes_{\bZ}\bR),
  \end{align*}
  then any $\alpha^{i}_{n,k}$ from above satisfies $\phi_{n}(\alpha^{2}_{n,k})=\alpha^{1}_{n,k}\otimes\alpha^{3}_{n,k}$.
\end{prp}


\section{estimation at infinity}

\begin{mtv}
  To estimate the volume, the key is to use the sequence
  \[
    \begin{tikzcd}
      0\ar[r] & \HH{}{0}(X,n\cL)\ar[r, "s"]\ar[r] & \HH{}{0}(X,(n+1)\cL)\ar[r] & \HH{}{0}(Y,(n+1)\cL|_{Y})\ar[r] & 0
    \end{tikzcd}
  \]
  in which $s\in\HH{}{0}(X,\cL)$ general nonzero section, $Y=\Divs(s)$.

\end{mtv}

\begin{dfn}
  Let $V_{X,L^{2}}^{n}=\Vol(\HH{}{0}(X,n\cL))$ and $V_{Y,L^{2}}^{n}=\Vol(\HH{}{0}(Y,n\cL|_{Y}))$ be the standard $L^{2}$-norm, define
  \[
    g(n+1)=\frac{V_{X,L^{2}}^{n+1}}{V_{X,L^{2}}^{n}\otimes V_{Y,L^{2}}^{n+1}}.
  \]

  Moreover, there is an $L^{2}$-norm on $\HH{}{0}(X,n\cL)$ induced by $s$, denote by $V_{s,L^{2}}^{n}$ the volume in this case; there is a quotient $L^{2}$-norm on $\HH{}{0}(Y,(n+1)\cL|_{Y})$, denote by $V_{q,L^{2}}^{n+1}$ in this case. Define
  \[
    V_{X,L^{2}}^{n}=\delta(n)V_{s,L^{2}}^{n},\quad V_{Y,L^{2}}^{n}=\gamma(n)V_{q,L^{2}}^{n},\quad\varphi(n+1)=\frac{V_{X,L^{2}}^{n+1}}{V_{s,L^{2}}^{n}\otimes V_{q,L^{2}}^{n+1}}.
  \]
\end{dfn}
\begin{rmk}
  So the next goal will be estimating these $g(n+1), \varphi(n+1), \delta(n), \gamma(n)$.
\end{rmk}


\begin{prp}\label{prp:estG}
  Let $P$ be the Hilbert polynomial of $\cL$ on $X$, then
  \[
    \lim_{n\ar\infty}\frac{\log g(n+1)}{P(n)}=-\int_{X}\log\sP{s(x)}^{2}\dif x.
  \]
\end{prp}

\begin{prp}
  Take an $L^{2}$-orthonormal basis $s_{1},\ldots,s_{N}, N=P(n)$ of $\HH{}{0}(X,n\cL)$, then $b_{n}(x)=\sum_{j=1}^{N}\sP{s_{j}(x)}^{2}$ is independent of the choice of basis and
  \[
    \lim_{n\ar\infty}\frac{b_{n}(x)}{P(n)}=1
  \]
  convergent uniformly in $x$.
\end{prp}
\begin{rmk}
  This result is off the shelf.
\end{rmk}


\begin{lem}
  \begin{enumerate}
    \item $\log\varphi(n)=o(n^{r})$,
    \item $\log\gamma(n)=o(n^{r})$,
    \item $\log\delta(n)=P(n)\int_{X}\log\sP{s(x)}^{2}\dif x + o(n^{r})$.
  \end{enumerate}
\end{lem}
\begin{rmk}
  These are for Proposition~\ref{prp:estG}. Almost all Chapter~3 is devoted to prove these. So Let's accept Proposition~\ref{prp:estG}.
\end{rmk}

\begin{prp}[Generalization of Proposition~\ref{prp:estG}]
  Let $\cL$ be ample, $j>0$ such that $j\cL$ very ample, fix $s\in\HH{}{0}(X,j\cL)$ such that $\Divs(s)=Y$ is smooth when $r\geq 2$ and reduced when $r=1$. Now for $n\gg 0, 0\leq p\leq j-1$, we have short exact sequence
  \[
    \begin{tikzcd}
      0\ar[r] & \HH{}{0}(X,(nj+p)\cL)\ar[r, "s"] & \HH{}{0}(X, ((n+1)j+p)\cL)\ar[r] & \HH{}{0}(Y, ((n+1)j+p)\cL|_{Y})\ar[r] & 0.
    \end{tikzcd}
  \]
  Define $g_{p}(n)$ via
  \[
    g_{p}(n+1)=\frac{V_{X,L^{2}}^{(n+1)j+p}}{V_{X,L^{2}}^{nj+p}\otimes V_{Y,L^{2}}^{(n+1)j+p}}.
  \]
  Analogously can show
  \[
    \lim_{n\ar\infty}\frac{\log g_{p}(n)}{P(nj)}=-\int_{X}\log\sP{s(x)}^{2}\dif x.
  \]
\end{prp}


\section{arithmetic Hilbert-Samuel theorem}

\begin{stp}
  Recall the notation: on $\HH{}{0}(X,n\cL)\otimes_{\bZ}\bR$, we have three norms $L^{2}$ norm gives volume $V_{X,L^{2}}^{n}$, the sup norm gives volume $V_{X, \mathrm{sup}}^{n}$, and the volume form $\alpha_{n,k}$ defined viewing $\HH{}{0}(X,n\cL)$ as a graded module (Definition~\ref{dfn:volk}) gives the volume $V_{X,k}^{n}$. Now compare them via
  \begin{align*}
    V_{X,k}^{n}&=f_{X}(n,k)V_{X,L^{2}}^{n},\\
    V_{X,L^{2}}^{n}&=h_{X}(n)V_{X, \mathrm{sup}}^{n}.
  \end{align*}
\end{stp}

\begin{prp}
  \begin{enumerate}
    \item $\log h_{X}(n)=o(n^{r+1})$,
    \item there is a unique real $k$ such that $\log f_{X}(n,k)=o(n^{r+1})$.
  \end{enumerate}
\end{prp}
\begin{rmk}
  The first is easy. The second is harder, need to discuss in very ample case and general case.

  With all these estimates we are ready for the two main theorems (Theorem~\ref{thm:supub}, Theorem~\ref{thm:L2ub}).
\end{rmk}


\begin{prf}[Theorem~\ref{thm:supub}, Theorem~\ref{thm:L2ub}]
  It suffices to use $P(n)=O(n^{r}), \sum_{j=0}^{n-1}P(j)=O(n^{r+1})$ (Faulhaber's formula), and that $\log(f_{X}(n,k)h_{X}(n))=o(n^{r+1}), \log h_{X}(n)=o(n^{r+1})$ above. The key is to relate the volume to Hilbert polynomial.
\end{prf}



\section{comparison with intersection theories}

Omitted.


% \subsection{case for arithmetic surface}

% \subsection{general case}



% ----------------------------------------------------------------------------------------------------
\bibliographystyle{alpha}
\bibliography{../universal-ref}
\end{document}
